% Generated draft: requires mathematical and source review.
\documentclass[11pt]{article}
\usepackage[margin=25mm]{geometry}
\usepackage{fontspec}
\setmainfont{Latin Modern Roman}
\usepackage{amsmath,amssymb,mathtools,unicode-math}
\setmathfont{Latin Modern Math}
\usepackage{enumitem,needspace,xurl,hyperref}
\setlength{\parindent}{0pt}
\setlength{\parskip}{5pt}
\setlength{\emergencystretch}{3em}
\newcommand{\N}{\mathbb{N}}
\newcommand{\Z}{\mathbb{Z}}
\newcommand{\Q}{\mathbb{Q}}
\newcommand{\R}{\mathbb{R}}
\newcommand{\C}{\mathbb{C}}
\newcommand{\F}{\mathbb{F}}
\newcommand{\A}{\mathbb{A}}
\newcommand{\PP}{\mathbb{P}}
\newcommand{\E}{\mathbb{E}}
\newcommand{\eps}{\varepsilon}
\newcommand{\dd}{\,\mathrm{d}}
\begin{document}
\section*{UnsolvedMath ID: 700098}
\textbf{Source problem code:} catalogue entry
\subsection{Definitions and mathematical statement}
An abstract group-theoretic property \(\Sigma\) is understood as a property invariant under group isomorphism. For every group \(G\), define
\[
\Sigma(G)=\left\langle\bigcup_{\substack{H\trianglelefteq G\\H\text{ has property }\Sigma}}H\right\rangle.
\]
Thus \(\Sigma(G)\) is the subgroup generated by all normal subgroups of \(G\) having property \(\Sigma\); the subgroup generated by an empty family is \(\{1\}\). This subgroup is normal in \(G\).

The property \(\Sigma\) is \emph{radical in the sense of this problem} if
\[
\text{for every group }G,\qquad \Sigma(G)\text{ has property }\Sigma.
\]
When this holds, \(\Sigma(G)\) is called the \(\Sigma\)-radical of \(G\). A radical property \(\Sigma\) is \emph{strongly radical} if, for every group \(G\),
\[
\text{for every }K\trianglelefteq G/\Sigma(G),\qquad
K\text{ has property }\Sigma\ \Longrightarrow\ K=\{1\}.
\]
The quantifiers here range over arbitrary groups, without a finiteness assumption.

\textbf{Part (a), the remaining question in the supplied record.} Is the property denoted in the source by \(\overline{\! RN}\) radical in this sense? Is it strongly radical in this sense? The supplied record does not define \(\overline{\! RN}\). Consequently, these questions specify the radical and strongly radical conditions precisely, but do not determine the class of groups to which those conditions are to be applied. No interpretation of this notation is imposed here.

\textbf{Part (b), archived as solved in the supplied record.} Is the property \(N^0\), referenced there to Archive, Problem 1.38, radical in the same sense? The supplied record does not include the definition of \(N^0\) or the text of Archive, Problem 1.38.

The supplied archive reports a negative answer to part (b), attributed to Yu. I. Merzlyakov (1973). Its reported counterexample is \(G=\operatorname{SL}_n(\Z)\), the group of determinant-one \(n\times n\) integer matrices, for sufficiently large integers \(n\geq 3\). The archive states that this group contains a nonabelian finite simple subgroup and therefore is not an \(N^0\)-group. It also states that
\[
G=\Gamma_n(2)\Gamma_n(3),\qquad
\Gamma_n(m)=\ker\bigl(\operatorname{SL}_n(\Z)\longrightarrow
\operatorname{SL}_n(\Z/m\Z)\bigr),\quad m\in\{2,3\},
\]
and that these congruence subgroups have central systems and hence are \(N^0\)-groups, citing Kargapolov and Merzlyakov and referring again to Archive, Problem 1.38. These are supplied historical assertions, not independently verified conclusions of this draft. The record gives neither an explicit threshold for \(n\) nor a definition of a central system.
\subsection{Short English statement}
The problem asks whether the group property \(\overline{\! RN}\) is radical and whether it is strongly radical: normal subgroups with the property should generate a subgroup with the property, and the quotient by that subgroup should have no nontrivial normal subgroup with the property. The meaning of \(\overline{\! RN}\) is missing from the supplied record. A separate question about \(N^0\) is archived with a reported negative answer; its definition is also missing. The supplied registry labels the overall problem partially solved.
\subsection{Sources}
\begin{itemize}
\item Local UnsolvedMath registry, supplied complete source record: permanent ID \(700098\), exact source code \texttt{catalogue entry}, title \emph{Kourovka Notebook Problem 2.22}. The supplied status and canonical status are \texttt{partially\_solved}. The status review is dated \texttt{2026-10-16}, records its evidence as \texttt{published}, and describes archived solved parts with remaining parts in the main section. This is supplied dated curation, not an independently verified mathematical or literature assessment.
\item \emph{The Kourovka Notebook (Unsolved Problems in Group Theory)}, 2nd issue (1967), Problem 2.22, proposed by Sh. S. Kemkhadze; updated in the 21st issue (2026), identified in the supplied record as \texttt{arXiv:1401.0300v48}. The supplied collection description names E. I. Khukhro and V. D. Mazurov as editors of the Issues 1--21 collection.
\item Supplied Archive of Solved Problems entry for Problem 2.22(b): negative answer and counterexample attributed to Yu. I. Merzlyakov (1973). The entry refers to Archive, Problem 1.38, for \(N^0\) and the central-system qualification; that referenced entry was not supplied.
\item M. I. Kargapolov and Yu. I. Merzlyakov, in: \emph{Itogi Nauki. Algebra. Topologiya. Geometriya, 1966}, VINITI, Moscow, 1968, pp. 57--90 (Russian), as cited in the supplied archive for the congruence-subgroup assertions. No article title is supplied.
\end{itemize}
\end{document}
